% ISBI full-paper submission -- ATDL assignment 1
% Build:  latexmk -pdf main.tex     (bibtex, NOT biber -- IEEEbib.bst is a
% BibTeX style)
% Requires in this directory: spconf.sty, IEEEbib.bst
% --------------------------------------------------------------------------
\documentclass{article}
\usepackage{spconf,amsmath,amssymb,graphicx}

% Times text + Times math, as the IEEE/ISBI guidelines ask for.
\usepackage{mathptmx}

% Compressed lists are explicitly allowed by the ISBI guidelines.
\usepackage{enumitem}
\setlist{nosep, leftmargin=14pt}

\usepackage{booktabs}     % proper rules in tables
% \usepackage{subcaption}   % (a)/(b) subfigures
\usepackage{microtype}    % subtle spacing tweaks; buys you ~half a line per
                          % column
\usepackage{mwe}         % supplies "example-image" placeholders -- DELETE
                         % once real figures are in
\graphicspath{{figures/}}

% hyperref last, cleveref after hyperref.
\usepackage[hidelinks]{hyperref}
\usepackage[capitalise,noabbrev]{cleveref}

% Handy macros -- keep notation consistent instead of retyping it.
\newcommand{\x}{\mathbf{x}}
\newcommand{\y}{\mathbf{y}}
\newcommand{\loss}{\mathcal{L}}
\newcommand{\model}{\textsc{OurMethod}}  % rename once, updates everywhere

% Title and authors
% -----------------
\title{On Accrediting Compression the Property of Generalization}

\name{Julian Barragan Gutierrez}
\address{Advanced Topics in Deep Learning \\
         University of Copenhagen}


\begin{document}

% Do NOT uncomment \ninept -- the assignment forbids changing the font size.
%\ninept

\maketitle

\begin{abstract}
	Despite their great success, we have yet to theoretically explain the
	generalization capabilities of Deep Neural Networks (DNNs).
	\cite{shwartzziv2017} attempt to give an information theoretical theory of
	deep learning, examining it through the lens of information theory, and
	claiming that the (emergent and somewhat inexplicable) generalization
	capabilities of DNNs are owed to compression occurring as a result of the
	stochasticity of the Stochastic Gradient Descent (SGD) optimization method.

	In a rebuttal paper \cite{saxe2018} the claims and methodology of
	aforementioned paper is largely disproven. In this paper we scrutinize
	the Information Bottleneck Theory of Deep Learning claims, the rebuttal
	papers validity, and focus our discussion on the accreditation of
	generalization to compression.
\end{abstract}

\begin{keywords}
	Deep Neural Networks, Information Bottleneck, compression, generalization
\end{keywords}

% !TEX root = ../main.tex
\section{Synopsis}\label{sec:synopsis}

\textit{The Information Bottleneck (IB) theory of deep learning} is first 
introduced
in \cite{tishby2015}, building on the IB principle of \cite{tishby1999}, and
examined empirically in \cite{shwartzziv2017}. Here we will focus our attention
to the claims and results presented in \cite{shwartzziv2017}. They introduce a
promising model invariant methodology for each hidden layer $T_i$ in a DNN,
where each hidden layer $T_i$ is treated as one random variable, so that the
network forms the Markov chain $Y\to X\to T_1\to\dots\to T_k\to\hat{Y}$.
This measures what they in
\cite{shwartzziv2017} interpret as the compression and generalization 
capabilities
by respectively measuring the mutual information $I(X;T)$ and $I(Y;T)$. This
creates coordinates in the information plane \cite{tishby2015}.
The mutual information can't be measured directly, so they solve this by
binning the activations w.r.t. the $2^{12}=4096$
different input patterns into 30 equal interval bins. This is inevitable since 
$T$ is in principle a deterministic function of $X$ making the mutual 
information $I(X; T)$ infinite.

\textit{To argue the SGD Drift-diffusion phase}
they examine the behaviour of DNNs on said information plane as a function of
number of epochs. They claim
this provides evidence that training is split into two main phases; a short
initial fitting phase reducing training error, and one much longer compression
phase induced by stochastic relaxation or random diffusion. Thus owed to SGD
learning transitions into what they call the \textit{SGD drift-diffusion phase}
whose transition they claim is traceable both in the information plane, as
well as by the weight gradients mean norm and standard deviation.

\textit{The loss function of choice} for networks trained in 
\cite{shwartzziv2017} is
the \textit{cross-entropy loss function}. This choice is explained simply by
stating that it is conventional for stochastic optimization. They specifically
mention that they expect minimization of the Empirical Error Minimization (ERM)
to be - reasonably so - a result of the cross-entropy loss function. Seemingly
no more consideration as to a potential influence in compression is mentioned.

The cross-entropy loss or log-loss is defined as
\begin{equation}
	CE = -\sum_{i=1}^{N} q_i\log_2(p_i)
\end{equation}

% !TEX root = ../main.tex
\section{What Instigates Compression in DNNs}\label{sec:compression}
A rebuttal paper released shortly after \cite{shwartzziv2017} written by Saxe
et al. \cite{saxe2018} provide evidence that several claims from the former
paper are wrong. Here we will focus our attentions to what has a direct causal
link to the observed compression phenomena as seen in the information plane, and
measured by the mutual information $I(X; T)$.

In \cite{shwartzziv2017} they claim that the stochasticity introduced by
SGD induces a random diffusion process which can be described by the
Fokker-Planck equation. Their argument is that this diffusion, subject to the
error constraint, maximizes the conditional entropy $H(X\mid T_i)$ i.e.
minimizes the mutual information $I(X;T_i)=H(X)-H(X\mid T_i)$, since $H(X)$ is
fixed by the data. Thus attributing the compression to the optimization method
SGD. In \cite{saxe2018} they show that compression is seen
as a result of:
\begin{itemize}
	\item \textbf{Double-saturating nonlinearities} cause the hidden activations
	to saturate as training progresses. The binned activity then collapses into
	the two extreme bins, and the measured $I(X;T)$ drops. Their tanh network
	compresses, their ReLU network does not. The same holds when the mutual
	information is estimated with kernel density and $k$-nearest-neighbour
	estimators instead of binning.
	\item \textbf{Growth in weight magnitude} make the units enter this
	saturated regime in the first place. In their minimal model of a single
	unit $h=f(w_1X)$ with $X\sim\mathcal{N}(0,1)$, $I(T;X)=H(T)$, since $T$ is
	deterministic given $X$, and they show $H(T)$
	first rises and then falls as a function of the weight $w_1$ for tanh,
	whilst it only rises for ReLU.
	\item \textbf{Genuinely task-irrelevant input dimensions} do compress, but
	this happens at the same time as the fitting and not in a separate phase
	afterwards. The information on the input as a whole can still increase.
\end{itemize}

% Your own saturation/binning sketch goes here. Uncomment once figures/ has it.
% \begin{figure}[t]
%   \centering
%   \includegraphics[width=\linewidth]{saturation}
%   \caption{}
%   \label{fig:saturation}
% \end{figure}


Subsequently they find that compression is \textit{not} a result of:
\begin{itemize}
	\item \textbf{The drift-diffusion transition} in the gradient SNR. They
	find this transition in all cases, also in networks that never compress,
	so the two phases are not causally related to the compression.
	\item \textbf{The maximum entropy argument} of \cite{shwartzziv2017}, which
	they argue does not hold theoretically. The stationary distribution it
	describes is a distribution over weights across different training runs,
	whilst $H(X\mid T)$ is about the inputs, and this gives no reason that the
	weights found in one particular run maximize $H(X\mid T)$. They do not show
	the claim to be false, but leave the burden of proof with
	\cite{shwartzziv2017}.
	\item \textbf{The stochasticity of SGD}, as they train the same networks
	with full batch gradient descent, which has no randomness in its updates,
	and still see comparable compression in the tanh networks.
\end{itemize}
where their account thus reduces to a single mechanism; the weights grow during
training, and the double-saturating nonlinearity turns this growth into a drop
in the measured $I(X;T)$. However the mechanism relies on the weights growing in
the first place, and what makes them grow is never asked, but simply taken as
given. The loss function is held fixed in all of the comparisons they make. The
permutation which would settle this, tanh trained with MSE, is never run, and
the only setting where they do use MSE is the linear network, in which the
nonlinearity is changed as well.

% !TEX root = ../main.tex
\section{Discussion}\label{sec:discussion}

\textit{The loss as the variable nobody turned}

\textit{Cross-entropy as a drive toward saturation}

\textit{What this account explains}

\textit{The missing permutation}


\bibliographystyle{IEEEbib}
\bibliography{refs}

\end{document}
